\section{Introduction}
\label{sec:introduction}

SPH is well known for its adaptivity; by tracing mass rather than volume the method can accurately capture the large dynamic range (of a factor $10^{11}$ in density, and $10^4$ in smoothing length, in cosmological galaxy formation simulations, see \cite{Borrow2018}) present in astrophysical phenomena, and is hence a natural choice for such simulations. The post-processing of SPH data in adaptive environments, however, presents a challenge due to the unpredictable structure of the underlying particle fields.

One post-processing case of interest, particularly within the astrophysics community, is the projection of three dimensional data onto a fixed two-dimensional grid. In astrophysics, projected quantities naturally correspond to those that are observable; for instance, it is significantly easier to measure the gas column density of a galaxy than it is to probe the three dimensional density structure. With projected quantities being instrumental in the connection of simulated datasets and the real Universe, obtaining accurate projected reconstructions of density (and other, e.g. temperature) fields from SPH data is vital. In specialised post-processing fields such as radiative transfer convergence and accuracy are a subject of regular discussion (e.g. \cite{Yang2020}), but for projected SPH quantities it is generally assumed that basic algorithms will suffice.

Typical solutions to the `projection problem' do not guarantee converged results in highly adaptive situations, particularly in the case where the smoothing length is of a similar order to the cell size that is chosen for projection. In this context, convergence refers to producing a result at a given grid resolution that corresponds to a resolution where all particles are resolved by a large number of cells, down-sampled to the lower resolution. Past work (e.g. \cite{Koepferl2017, Petkova2018}) introduced the concept of using a Voronoi decomposition to ensure that each pixel is assigned the correct fraction of each particle. This procedure, however, is computationally and conceptually complex.

We consider various methods to project SPH data to a fixed grid, from basic kernel interpolation, to `blitting' kernels to the fixed grid, and subsampling the kernel. There are many software packages available to perform this deposition, including (but not limited to) those described in \cite{Price2007, Turk2011, Benitez-Llambay2015}. The open-source package used in this work is {\tt swiftsimio} \cite{Borrow2020a}, used for reading and visualising data produced by the {\tt SWIFT} \cite{Schaller2016} code, and is the first package that we are aware of that makes subsampled techniques publicly available.
